\section{Discussion}
\label{Sec:Discussion}

The model’s strong performance metrics confirm the efficacy of autoencoders architectures for anomaly detection tasks. It shows that the model generalizes well from training exclusively on normal data, successfully detecting unseen anomalies during testing. The last is key for developing cost-effective solutions as the collection of anomalous samples may be considerably expensive.

The results also highlight the importance of threshold selection. An F1-optimized threshold yielded a balanced trade-off between precision and recall. However, for industrial applications other approaches may be explored, specially those taking into account the difference between the cost of a false alarm (false positive) and a non identified anomaly (false negative). 

Despite these promising results, one must note the limitation due to the reduced experiment's scope. The evaluation was focus to the Slide Rail machine type so generalization to other machine types is not guaranteed.